\documentclass[11pt]{article}
\usepackage[margin=1in]{geometry}
\usepackage{amsmath,amssymb}
\usepackage{hyperref}

\title{A Simple Accessibility Theory of Parity Circuit Competition}
\date{}

\begin{document}
\maketitle

\section{Setup}

Consider a pure-ReLU residual block. Suppose a neuron combines \(p\) orthogonal parent Walsh modes
\[
q_1,\ldots,q_p,\qquad \|q_i\|=Q,
\]
with balanced signs. If \(v\) is the neuron's weight vector, the preactivation amplitude of each parent is
\[
\gamma=\frac{\|v\|Q}{\sqrt p}.
\]

For a balanced \(p\)-parent ReLU gate, the coefficient of the product Walsh mode is
\[
c_p=\kappa_p\gamma,
\]
with
\[
|\kappa_2|=\frac12,\qquad
|\kappa_4|=\frac14,\qquad
|\kappa_8|=\frac{5}{32}.
\]

Define the accessibility efficiency
\[
R\equiv \frac{|c_p|}{\|v\|}
=\frac{|\kappa_p|Q}{\sqrt p}.
\]
Thus \(R\) measures how much of the desired child mode is produced per unit hidden-weight norm. The simple theory predicts that, when several routes to the same target mode are simultaneously available, the route with larger \(R\) is more accessible.

\section{\texorpdfstring{\(2+2\to4\)}{2+2 to 4} versus \texorpdfstring{\(1+1+1+1\to4\)}{1+1+1+1 to 4}}

Let \(Q_1\) be the norm of each relevant singleton mode and \(Q_2\) the norm of each relevant degree-2 mode.

For the direct route \(1+1+1+1\to4\), we have \(p=4\), \(Q=Q_1\), and \(|\kappa_4|=1/4\):
\[
R_{1^4}
=
\frac{(1/4)Q_1}{\sqrt4}
=
\frac{Q_1}{8}.
\]

For the compositional route \(2+2\to4\), we have \(p=2\), \(Q=Q_2\), and \(|\kappa_2|=1/2\):
\[
R_{2+2}
=
\frac{(1/2)Q_2}{\sqrt2}
=
\frac{Q_2}{2\sqrt2}.
\]

The two routes are equally efficient when
\[
\frac{Q_2}{2\sqrt2}
=
\frac{Q_1}{8},
\]
hence
\[
\boxed{
\frac{Q_2}{Q_1}
=
\frac{1}{\sqrt8}
\approx 0.3536.
}
\]

Therefore
\[
\frac{Q_2}{Q_1}>\frac1{\sqrt8}
\]
predicts that \(2+2\to4\) is more accessible, while below this ratio the direct \(1^4\to4\) route is more accessible.

\section{\texorpdfstring{\(4+4\to8\)}{4+4 to 8} versus \texorpdfstring{\(1+\cdots+1\to8\)}{direct degree-8}}

Let \(Q_4\) be the norm of each relevant degree-4 parent mode.

For the direct route \(1+\cdots+1\to8\), we have \(p=8\), \(Q=Q_1\), and \(|\kappa_8|=5/32\):
\[
R_{1^8}
=
\frac{5}{32}\frac{Q_1}{\sqrt8}.
\]

For \(4+4\to8\), we have \(p=2\), \(Q=Q_4\), and \(|\kappa_2|=1/2\):
\[
R_{4+4}
=
\frac{Q_4}{2\sqrt2}.
\]

Equating the two,
\[
\frac{Q_4}{2\sqrt2}
=
\frac{5Q_1}{32\sqrt8}.
\]
Since \(\sqrt8=2\sqrt2\),
\[
\boxed{
\frac{Q_4}{Q_1}
=
\frac{5}{32}
=
0.15625.
}
\]

Therefore
\[
\frac{Q_4}{Q_1}>0.15625
\]
predicts that \(4+4\to8\) is more accessible, while below this ratio the direct \(1^8\to8\) route is more accessible.

\section{Experimental method}

The theory can be tested while separating circuit accessibility from the dynamics by which the parent features are learned.

\begin{enumerate}
    \item Train a first ReLU residual block to produce the required parent parity modes. For the degree-4 experiment, train degree-2 parents; for the degree-8 experiment, train degree-4 parents.

    \item Freeze the embedding and the trained first block.

    \item Rescale the frozen block by
    \[
    V^{(1)}\mapsto \alpha V^{(1)}.
    \]
    By ReLU homogeneity,
    \[
    \operatorname{ReLU}(\alpha z)=\alpha\operatorname{ReLU}(z),
    \]
    so this changes the strength of the learned parent modes without changing the learned circuit itself.

    \item For every \(\alpha\), directly measure the Walsh-mode norms in the frozen residual stream, in particular
    \[
    \frac{Q_2}{Q_1}
    \quad\text{or}\quad
    \frac{Q_4}{Q_1}.
    \]
    These measured ratios, rather than \(\alpha\) itself, are the theoretical control variables.

    \item Add a fresh ReLU block and readout and train only them on the higher-degree target.

    \item Determine which circuit forms from the Walsh decomposition of the new block's preactivations. A \(2+2\to4\) neuron should be dominated by its two degree-2 parents; a \(4+4\to8\) neuron by its two degree-4 parents; a direct route should be dominated by the relevant singleton modes.
\end{enumerate}

The key predictions are
\[
\boxed{\frac{Q_2}{Q_1}=\frac1{\sqrt8}}
\]
for the degree-4 experiment and
\[
\boxed{\frac{Q_4}{Q_1}=\frac5{32}}
\]
for the degree-8 experiment.

These predictions assume pure ReLU, approximately equal-norm and orthogonal parent modes, and balanced gates.

\end{document}
